\documentclass[10pt,a4paper]{amsart}
\usepackage[margin=19mm]{geometry}
\usepackage{amsmath,amssymb,mathtools}
\usepackage[scr=rsfs,cal=euler]{mathalfa}
\usepackage{xfrac}
\usepackage[unicode,hidelinks,bookmarks=false]{hyperref}
\newcommand{\ie}{i.e.\ }
\newcommand{\cf}{cf.\ }
\newcommand{\resp}{respectively\ }
\newcommand{\Subsection}{\subsection}
\providecommand{\nobreakdash}{\nobreak\hbox{-}}
\setlength{\emergencystretch}{2em}
\multlinegap0pt
\usepackage{bm}
\usepackage{bbm}
\usepackage{dsfont}
\usepackage{xspace}
\usepackage{cancel}
\usepackage{mathtools}
\usepackage{amsthm}
\makeatletter
\@ifundefined{lemma}{\newtheorem{lemma}{Lemma}}{}
\@ifundefined{theorem}{\newtheorem{theorem}[lemma]{Theorem}}{}
\makeatother
\DeclareMathOperator{\supp}{supp}
\newcommand{\Rd}{{\R^d}}
\newcommand{\R}{\mathbb{R}}
\newcommand{\uinf}{u_\infty}
\newcommand{\ximin}{\xi_{\text{min}}}
\title[Nonlinear Diffusion Equations: Full characterization of Entr]{Nonlinear Diffusion Equations: Full characterization of Entropies}
\author{Anton Arnold, Jose A. Carrillo, Daniel Matthes}
\date{Source: arXiv:2608.07129v1 (CC BY 4.0)}
\begin{document}
\maketitle

\noindent\emph{Source and licence.} Nonlinear Diffusion Equations: Full characterization of Entropies. Version arXiv:2608.07129v1 (CC BY 4.0), distributed under the Creative Commons Attribution 4.0 International licence, as recorded in the arXiv licence metadata for this exact version and shown on the abstract page https://arxiv.org/abs/2608.07129v1. Full rights notice: licenses/arxiv-2608.07129-CC-BY-4.0.txt.\par\smallskip

\noindent\textit{Editorial scope.} Counted unit: the Theorem whose source label is \texttt{gencz} in the article (source0.tex lines 2778--2786, source label \texttt{gencz}), delivered together with the article's own complete original proof of it (source0.tex lines 2788--2925, one \texttt{proof} environment of 138 lines). No mathematics is written, repaired or re-derived: statement and proof are the article's own text line for line. Required context retained uncounted: Taylorent (lines 436--438). Every definition, display and label of those lines is retained unchanged. Omitted from this package and not redistributed: 1--435, 439--2777, 2787--2787, 2926--3606. Nothing inside a retained excerpt is omitted or altered: every retained body is delivered exactly as the article's lines are written, so no omission receipt of any kind accompanies this package. Citations: the retained excerpts carry no citation command at all, so no bibliography entry of the article is transcribed and no reference is redistributed. Reference adaptation: none was needed and none was made; every reference inside the delivered unit resolves inside it, so no unresolved reference or citation is produced. Printed slips kept literal: no printed word, symbol, hypothesis or number of the counted statement or of its proof was altered. Layout and class: the article's own class is \texttt{amsart} with packages including amsmath, amsthm, amssymb, amsfonts, yfonts, mathtools, hyperref, color, cases, enumitem, empheq, graphicx, multimedia, longtable; the wrapper is the lane's pinned \texttt{amsart} editorial wrapper at 10pt on A4, which restates the theorem-like environments the retained text needs and adds the article's own macro definitions that the retained text needs (\texttt{\textbackslash R}, \texttt{\textbackslash Rd}, \texttt{\textbackslash supp}, \texttt{\textbackslash uinf}, \texttt{\textbackslash ximin}), with extra wrapper packages (bm, bbm, dsfont, xspace, cancel, mathtools). The delivered numbering is the unit's own. Assets: no retained span contains a figure environment, a graphics-inclusion command, a PDF-page inclusion, a table or a TikZ picture; no image file is copied into this package, the package declares no asset dependency and no image file is redistributed with it. Licence: version arXiv:2608.07129v1 of this article is distributed under the Creative Commons Attribution 4.0 International licence, as recorded in the arXiv licence metadata for this exact version and shown on the abstract page https://arxiv.org/abs/2608.07129v1. A full rights notice is delivered at licenses/arxiv-2608.07129-CC-BY-4.0.txt. Characters: every retained excerpt is pure ASCII, so the delivered engine, XeLaTeX with its default Latin Modern text font, reports no missing character.
\begin{equation}\label{Taylorent}
 \mathcal{H}_g(u|u_\infty) = \frac12 \int_{\R^d} g'(\phi(\eta(x))-\phi(\uinf(x))) \phi'(\eta(x)) (u(x)-\uinf(x))^2 \,dx\geq 0\,,
\end{equation}

\begin{theorem}\label{gencz}
Assume that either $P$ is a regular non-degenerate diffusion satisfying $\inf_{r\ge0} P'(r)>0$ or a degenerate diffusion in which case we further assume that for some constant $n>0$, the map $r\mapsto P'(r)/r^n$ is continuous and positive on $(0,\infty)$, and is non-increasing on some interval $(0,r_0]$. Moreover, assume that the entropy function $g$ is such that $g'$ is positive on $(\xi_{\min},\infty)$.
Then there is a constant $C$ (depending on $P$, $V$, $M$, $g$, and possibly $n$ for degenerate diffusions) such that
    \begin{align}
        \label{eq:ourCK}
        \|u-\uinf\|_{L^1} \le C {\mathcal H}_g(u|\uinf)^{1/2}
    \end{align}
for all functions $u\in L^1_M(\R^d)$ such that ${\mathcal H}_g(u|\uinf)<\infty$.\\
\end{theorem}

\begin{proof}
The proof needs to improve over the Taylor expansion in \eqref{Taylorent}. We structure it in the following steps:

\

{\bf Step 1.-}
In the regular non-degenerate diffusion case, we obtain
at every $x$ with $0\leq u(x)<\uinf(x)$ by integrating by parts using that $g(0)=0$:
 \begin{align}
        G\big(u,u_\infty\big)(x)
        &= \int^{u(x)}_{\uinf(x)} g\big(\phi(s)-\phi(\uinf(x))\big) \, ds \nonumber\\
        &= - (\uinf(x)-u(x))\,g(0) +\int
        _{u(x)}^{\uinf(x)}(s-u(x))\,g'\big(\phi(s)-\phi(\uinf(x))\big)\phi'(s)ds  \nonumber \\
        &\ge \int_{(\uinf(x)+u(x))/2}^{u_\infty(x)}(s-u(x))\,g'\big(\phi(s)-\phi(\uinf(x))\big)\phi'(s)ds \nonumber\\
        &\ge {\frac38} \left(\uinf(x)-u(x)\right)^2\inf_{\frac{\uinf(x)}2<s<\uinf(x)} g'\big(\phi(s)-\phi(\uinf(x))\big)\: \inf_{\frac{\uinf(x)}2<s<\uinf(x)}\phi'(s)\nonumber \\
        &=: {\frac38}\left(\uinf(x)-u(x)\right)^2 \, I_1\, I_2\,, \label{G-est-nondeg}
    \end{align}
where $I_1$ and $I_2$ are the two infima in the penultimate line.

In the degenerate diffusion case, recall that $\phi(u_\infty(x))+V(x)=\bar C$ for all $x\in B_\infty$. It then follows at every $x$ with $0\leq u(x) < \uinf(x)$ --- hence such an $x$ is in $B_\infty$ --- that
    \begin{align*}
        \tilde G\big(u(x),\uinf(x);x\big)
        &= \int^{u(x)}_{\uinf(x)} 
        g\big(\phi(s)+V(x)-\bar C\big)ds  \\
        &= \int^{u(x)}_{\uinf(x)} g\big(\phi(s)-\phi(\uinf(x))\big) \, ds 
        \geq {\frac38}\left(\uinf(x)-u(x)\right)^2 \, I_1\, I_2\,,
    \end{align*}
    proceeding similarly as for \eqref{G-est-nondeg}.

\

{\bf Step 2.-} The first infimum $I_1$ is positive, uniformly in $x\in B_\infty$. To see this, define $U:=\max u_\infty>0$, implying $\phi(U)=\bar C$. In both cases (regular non-degenerate diffusion and degenerate diffusion), $P'$ is continuous on $[0,U]$ and $\phi$ is monotonically increasing.    
In the degenerate case we recall that $\xi_{\min}=\phi(0)-\bar C=-\int_0^U P'(r)/r\,dr>-\infty$.
Hence we estimate for all $s\in(\uinf(x)/2,\uinf(x))$: 
\begin{align}\label{d-phi-estimate}
        0 &\ge \phi(s) - \phi(\uinf(x)) = - \int_s^{\uinf(x)}\frac{P'(r)\,dr}r 
        \ge - \int_{\uinf(x)/2}^{\uinf(x)}\frac{P'(r)\,dr}{r} \\
    &\ge
    \begin{cases}
        - \int_{0}^{U/2}\frac{P'(r)\,dr}{r}, \quad & u_\infty(x)\in[0,U/2]\,,\\
        - \int_{U/4}^{U}\frac{P'(r)\,dr}{r}, \quad & u_\infty(x)\in(U/2,U]
    \end{cases}  \nonumber \\
    &\ge \delta - \int_0^U \frac{P'(r)\,dr}{r}  = \delta+\xi_{\min} , \nonumber 
\end{align}
where $\delta:=\min\Big(\int_{0}^{U/4}\frac{P'(r)\,dr}{r},\,\int_{U/2}^{U}\frac{P'(r)\,dr}{r}\Big)>0$ is independent of $x\in B_\infty$. 
In the non-degenerate case we estimate the last integral in \eqref{d-phi-estimate} as:
$$
  - \int_{\uinf(x)/2}^{\uinf(x)}\frac{P'(r)\,dr}{r} 
        \ge -\|P'\|_{L^\infty(0,\infty)} \int_{\uinf(x)/2}^{\uinf(x)}\frac{dr}{r} =-\|P'\|_{L^\infty(0,\infty)} \log 2.
$$
In the degenerate case $g'$ is continuous and positive on $[\delta+\xi_{\min},0]$, and in the non-degenerate case on $[-\|P'\|_{L^\infty(0,\infty)} \log 2,0]$. Hence, in both cases there is a positive lower bound on $g'(\phi(s)-\phi(\uinf(x)))$ for all $s\in(\uinf(x)/2,\uinf(x))$. 

\

{\bf Step 3.-}
The second infimum $I_2$ is also positive: notice first that in the case of regular non-degenerate diffusions, we have
    \begin{align}\label{phi-est}
        \phi'(s) = \frac{P'(s)}{s} 
        \ge \frac{P'_\text{min}}{\uinf(x)} 
        \geq \frac{P'_\text{min}}{P'_{\text{max}}}  \frac{P'(\uinf(x))}{\uinf(x)}:= c\phi'(\uinf(x))
    \end{align}
for all $s\in(\uinf(x)/2,\uinf(x))$ with 
\[
    0<P'_\text{min} :=\min_{0\leq s\leq U} P'(s) \le P'_{\text{max}}:=\max_{0\leq s\leq U} P'(s). 
\]
By the assumed uniform positivity and continuity of $P'(r)/r^n$, 
we conclude that there is some $\bar c>0$ such that 
\[ 
    \frac{P'(s)}{s^n} \ge \bar c \frac{P'(r)}{r^n} 
\]
for all $0<s<r\le U$. Consequently,  we have for all $s\in(\uinf(x)/2,\uinf(x))$:
    \begin{align*}
        \phi'(s) = \frac{P'(s)}{s} 
        = s^{n-1}\frac{P'(s)}{s^n} 
        \ge \bar c s^{n-1}\frac{P'(\uinf(x))}{\uinf(x)^n}
        = \bar c\left(\frac{s}{\uinf(x)}\right)^{n-1}\frac{P'(\uinf(x))}{\uinf(x)}
        \ge \tilde c \phi'(\uinf(x))\,
    \end{align*}
with $\tilde{c}=\bar c\min(1,2^{1-n})$ achieving a similar estimate as in \eqref{phi-est} for the regular non-degenerate case.

\

{\bf Step 4.-} We now conclude by collecting the estimates of $I_1$ and $I_2$, that there is a constant $K$, independent of $u$, such that
 \begin{align*}
        \phi'(\uinf(x))(\uinf(x)-u(x))^2 \le K^2  G(u(x),\uinf(x))
    \end{align*}
in the regular non-degenerate diffusion case, and 
    \begin{align*}
        \phi'(\uinf(x))(\uinf(x)-u(x))^2 \le K^2 \tilde G(u(x),\uinf(x);x)
    \end{align*}
in the degenerate diffusion case, 
holds for all $x\in\R^d$ with $0\leq u(x)<\uinf(x)$. Now, both for non-degenerate and degenerate diffusion, recalling that $u$ and $u_\infty$ are non-negative and of the same mass $M$:
    \begin{align*}
        \|u-\uinf\|_{L^1} 
        &= 2\int_{\{u<\uinf\}} (\uinf-u)(x)\;dx \\
        & \le 2\left(\int_{\{u<\uinf\}}\phi'(\uinf)(u-\uinf)^2(x)\;dx\right)^{1/2}\left(\int_{u<\uinf}\frac{dx}{\phi'(\uinf)}\right)^{1/2} \\
        &\le 2K\left(\int_{B_\infty}\frac{dx}{\phi'(\uinf)}\right)^{1/2}\mathcal H_g(u|\uinf)^{1/2}
    \end{align*} 
leading to the desired estimate.

\
    
{\bf Step 5.-} It remains to show that 
    $$
    \int_{B_\infty}\frac{dx}{\phi'(\uinf)}<\infty\,.
    $$
In the regular non-degenerate diffusion case, notice first that $\supp(u_\infty)$ is $\R^d$ and we directly estimate the above integral as
    \begin{align*}
        \int_{\R^d}\frac{dx}{\phi'(u_\infty)}
        &= \int_{\R^d}\frac{\uinf}{P'(u_\infty)}dx\leq \frac1{P'_{\text{min}}}\int_{\R^d}\uinf dx=\frac{M}{P'_{\text{min}}}.
    \end{align*}
    In the degenerate diffusion case, define $\bar V:=\bar C-\phi(0+)=-\ximin$, so that $B_\infty=\{V\le\bar V\}$, and let $\rho:=\phi^{-1}(\bar C-\bar V/2)>0$. We have:
    \begin{align*}
        \int_{B_\infty}\frac{dx}{\phi'(u_\infty(x))}
        &= \int_{\{0<V<\bar V/2\}}\frac{dx}{\phi'\big(\phi^{-1}(\bar C-V(x))\big)}
        +\int_{\{\bar V/2< V<\bar V\}}\frac{dx}{\phi'\big(\phi^{-1}(\bar C-V(x))\big)}\\
        &\le \frac{\big|\{0<V<\bar V/2\}\big|}{\inf_{\rho<r< U}\phi'(r)}
        +\int_{\{\bar V/2< V<\bar V \}}\frac{dx}{\phi'\big(\phi^{-1}(\bar C-V(x))\big)}.
    \end{align*}
    Since $\phi'(r)=P'(r)/r$ has a positive lower bound on $[\rho, U]$ by hypothesis, the first expression is finite, and it suffices to estimate the remaining integral.

    For evaluation of the second integral, we pass to radial coordinates $x=x_0+r\theta$, where $\theta\in{\mathbb S}^{d-1}$ and $r\ge0$, and $x_0$ is the minimal point of $V$. Let $0<\underline r(\theta)<\overline r(\theta)$ such that
    \[ \{x\in\Rd\ |\ \bar V/2< V(x) < \bar V\} = \{ x_0+r\theta\ |\ \theta\in\mathbb{S}^{d-1},\ \underline r(\theta)<r<\overline r(\theta) \}. \]
    Since $r\mapsto V(x_0+r\theta)$ is a $\lambda$-convex function for each fixed $\theta$, with minimum zero at $r=0$, we have that $V(x_0+r\theta)\ge \lambda/2\,r^2$, and in particular we have that $\overline r(\theta)\le R:=(2\bar V/\lambda)^{1/2}$. Since further $\partial_r V(x_0+r\theta)\ge\lambda r$ for each $r>0$, and
    \begin{align*}
        \partial_r\phi^{-1}(\bar C-V(x_0+r\theta)) = \frac{-\partial_rV(x_0+r\theta)}{\phi'\big(\phi^{-1}(\bar C-V(x_0+r\theta))\big)},
    \end{align*}
    we conclude that
    \begin{align*}
        \int_{\{ \bar V/2< V<\bar V\}}\frac{dx}{\phi'\big(\phi^{-1}(\bar C-V(x))\big)}
        &= \int_{{\mathbb S}^{d-1}} \int_{\underline r(\theta)}^{\overline r(\theta)} \frac{r^{d-1} dr}{\phi'\big(\phi^{-1}(\bar C-V(x_0+r\theta))\big)} d\theta\\
        &\le -\frac1\lambda\int_{{\mathbb S}^{d-1}} \int_{\underline r(\theta)}^{\overline r(\theta)} r^{d-2}\partial_r\phi^{-1}(\bar C-V(x_0+r\theta)) dr d\theta \\
        &\le R^{d-2}|\mathbb{S}^{d-1}|\,\big[\phi^{-1}(\zeta)\big]_{\zeta=\phi(0+)}^{\zeta=\phi(\rho)}
        \le (2\bar V/\lambda)^{(d-2)/2}|\mathbb{S}^{d-1}|\,\rho,
    \end{align*}     
    where we have used that $\bar C-V(x_0+\overline r(\theta)\theta)=\bar C-\bar V=\phi(0+)$ by definition of $\bar V$, and $\bar C-V(x_0+\underline r(\theta)\theta)=\bar C-\bar V/2=\phi(\rho)$ by definition of $\rho$. 
    Thus, the integral is finite.
\end{proof}

\end{document}
